% Generated draft: requires mathematical and source review.
\documentclass[11pt]{article}
\usepackage[margin=25mm]{geometry}
\usepackage{fontspec}
\setmainfont{Latin Modern Roman}
\usepackage{amsmath,amssymb,mathtools,unicode-math}
\setmathfont{Latin Modern Math}
\usepackage{enumitem,needspace,xurl,hyperref}
\setlength{\parindent}{0pt}
\setlength{\parskip}{5pt}
\setlength{\emergencystretch}{3em}
\newcommand{\N}{\mathbb{N}}
\newcommand{\Z}{\mathbb{Z}}
\newcommand{\Q}{\mathbb{Q}}
\newcommand{\R}{\mathbb{R}}
\newcommand{\C}{\mathbb{C}}
\newcommand{\F}{\mathbb{F}}
\newcommand{\A}{\mathbb{A}}
\newcommand{\PP}{\mathbb{P}}
\newcommand{\E}{\mathbb{E}}
\newcommand{\eps}{\varepsilon}
\newcommand{\dd}{\,\mathrm{d}}
\begin{document}
\section*{UnsolvedMath ID: 701487}
\textbf{Source problem code:} catalogue entry
\subsection{Definitions and mathematical statement}
An involution in a group is an element of order exactly \(2\). Let \(G\) be a group admitting a finite generating set \(\{s_1,\ldots,s_n\}\) consisting of involutions. Assume that
\[
(uv)^4=1_G
\qquad\text{for every pair of involutions }u,v\in G.
\]
The condition concerns all involutions of \(G\), not merely the chosen generators.

The source asks two questions:
\begin{enumerate}
\item Must \(G\) be finite?
\item Must \(G\) be a \(2\)-group?
\end{enumerate}
For possibly infinite groups, interpreting a \(2\)-group as a torsion \(2\)-group makes the second question precise as follows: must
\[
\forall g\in G\;\exists k\in\Z_{\geq 0}\quad g^{2^k}=1_G?
\]
The source does not explicitly specify its convention for the term \(2\)-group when finiteness is not known. For finite groups, this condition is equivalent to the group order being a power of \(2\).

The source states that the answers are affirmative for \(n\leq 3\). Its editors' comment of 2021 states that a previously published claim of an affirmative solution proved erroneous. These are supplied source qualifications; this draft does not independently verify them.
\subsection{Short English statement}
D. V. Lytkina asks whether every finitely generated group generated by involutions, in which the product of any two involutions has fourth power equal to the identity, must be finite and must be a \(2\)-group. The source states an affirmative answer for at most three generators and records a 2021 correction withdrawing a claimed general solution.
\subsection{Sources}
\begin{itemize}
\item Local UnsolvedMath registry, supplied record with permanent ID \(701487\), exact source problem code \texttt{catalogue entry}, title \emph{Kourovka Notebook Problem 18.58}. The supplied registry labels the problem open; that label is not an independently verified literature assessment.
\item \emph{The Kourovka Notebook (Unsolved Problems in Group Theory)}, Problem 18.58, proposed by D. V. Lytkina in 2014. The supplied background identifies the 18th issue (2014) and the 21st issue edition (2026), arXiv identifier \texttt{1401.0300v48}. The supplied collection description names E. I. Khukhro and V. D. Mazurov as editors.
\item Editors' comment of 2021, as reproduced in the supplied statement: a previously published claim of an affirmative solution proved erroneous. No bibliographic details for that claim are supplied.
\end{itemize}
\end{document}
